% Subproblem 1: Standardize the cube and write the T_{2g} mode explicitly
%
% CONTEXT: This is part of the proof that the Hessian of J_8 is positive definite
% on T_{2g}. We work in /data/projects/3eacb340-027e-4b71-8de4-67574ba54ea2/QBLYHzbdEjWeyg/proof/.
%
% NOTATION:
%  - (theta, phi) are standard spherical coordinates on S^2.
%  - G_q(z) = -(1/(2*pi)) * log(sin(dist(z,q)/2)) is the Green function on S^2 at q.
%  - For a point z in S^2, we write z = (sin(theta)cos(phi), sin(theta)sin(phi), cos(theta))
%    in Cartesian coordinates.
%
% STATEMENT TO PROVE:
%
% Let psi = phi - pi/4, and set x = sin(theta)cos(psi), y = sin(theta)sin(psi),
% z_coord = cos(theta), so that (x, y, z_coord) are the Cartesian coordinates of z.
%
% (a) Show that the eight cube vertices (in (theta, phi) coordinates)
%       p_1 = (theta_0, 0),          p_5 = (pi-theta_0, 0),
%       p_2 = (theta_0, pi/2),       p_6 = (pi-theta_0, pi/2),
%       p_3 = (theta_0, pi),         p_7 = (pi-theta_0, pi),
%       p_4 = (theta_0, 3*pi/2),     p_8 = (pi-theta_0, 3*pi/2),
%     (where theta_0 = arccos(1/sqrt(3))) are exactly the eight points
%     q_{eps1, eps2, eps3} = (eps1, eps2, eps3)/sqrt(3) in S^2,
%     for eps_j in {+1, -1}.
%
% (b) Show that in (theta, psi) coordinates, the sets
%       A = {p_1, p_2, p_7, p_8} and B = {p_3, p_4, p_5, p_6}
%     (which are the index sets for the T_{2g} generator: v_i = +partial_phi/sin(theta)
%      for i in A, v_i = -partial_phi/sin(theta) for i in B) correspond to
%       A = { q_eps : eps_1 * eps_3 = +1 }  and  B = { q_eps : eps_1 * eps_3 = -1 }.
%
% (c) Using the identity sin^2(dist(z,q)/2) = (1 - z.q)/2, show that
%       G_q(z) = -(1/(4*pi)) * log(1 - z.q) + const,
%     and derive the explicit formula
%       F_A(theta, psi) := partial_psi ( sum_{q in A} G_q(theta, psi) )
%         = (1/(4*pi*sqrt(3))) * sum_{eps: eps_1*eps_3=1}
%             (eps_2 * x - eps_1 * y) / (1 - (eps_1*x + eps_2*y + eps_3*z_coord)/sqrt(3)),
%     where the sum runs over (eps_1, eps_2, eps_3) in {+/-1}^3 with eps_1*eps_3 = 1.
%
% (d) Define F_B(theta, psi) = partial_psi (sum_{q in B} G_q) and show F_B(theta,psi)
%     = F_A(theta, psi+pi).
%
% (e) Define the symmetric and antisymmetric parts:
%       C(theta, psi) = F_A(theta, psi) + F_A(theta, psi+pi)
%         = (1/(4*pi*sqrt(3))) * sum_{all eps in {+/-1}^3}
%             (eps_2 * x - eps_1 * y) / (1 - (eps_1*x + eps_2*y + eps_3*z_coord)/sqrt(3)),
%       D(theta, psi) = F_A(theta, psi) - F_A(theta, psi+pi)
%         = (1/(4*pi*sqrt(3))) * sum_{all eps in {+/-1}^3}
%             eps_1*eps_3 * (eps_2*x - eps_1*y) / (1 - (eps_1*x+eps_2*y+eps_3*z_coord)/sqrt(3)).
%
% DELIVERABLE: A clean LaTeX proof of (a)-(e). Use the identity
%   partial_psi G_q = partial_psi [ -(1/(4*pi)) log(1 - (eps_1 x + eps_2 y + eps_3 z)/sqrt(3)) ]
%   = (1/(4*pi)) * (eps_1 partial_psi x + eps_2 partial_psi y) / (1 - ...)
% and note that partial_psi x = -y, partial_psi y = x.
%
% Output should be a self-contained LaTeX section (no \documentclass, just the content)
% suitable for \input{} into the main paper. Use the macros already defined in problem.tex:
% \uS, \sph, \dist, \R, etc.
\end{document}
